

\TypeQuestion{}{line_application_parallel}
\TextTall{図のように、放物線 $y=\dfrac{1}{2}x^2$ と直線 $\ell$ が2点 $\mathrm{A},\mathrm{B}$ で交わり、点 $\mathrm{A},\mathrm{B}$ の $x$ 座標はそれぞれ $-2,4$ である。直線 $m$ は直線 $\ell$ と平行で、放物線と2点 $\mathrm{C},\mathrm{D}$ で交わっている。点 $\mathrm{C}$ の座標を $(-4,p)$ とし、点 $\mathrm{D}$ の $x$ 座標は正である。次の問いに答えなさい。}
\QQ{点 $\mathrm{A},\mathrm{B}$ の座標をそれぞれ求めなさい。}{}
\QQ{直線 $\ell$ の式を求めなさい。}{}
\QQ{$x$ の値が点 $\mathrm{A}$ の $x$ 座標から点 $\mathrm{B}$ の $x$ 座標まで増加するとき、放物線における変化の割合を求めなさい。}{}
\QQ{$p$ の値と、直線 $m$ の式を求めなさい。}{}
\QQ{点 $\mathrm{D}$ の座標を求めなさい。}{}


\TypeQuestion{}{line_application_parameter}
\Text{図のように、放物線 $y=x^2$ 上に点 $\mathrm{A}(-2,4)$ と、$x$ 座標が正である点 $\mathrm{B}$ がある。点 $\mathrm{B}$ の $x$ 座標を $t$ とし、2点 $\mathrm{A},\mathrm{B}$ を通る直線を $\ell$ とする。また、直線 $\ell$ と $y$ 軸との交点を $\mathrm{C}$、原点を $\mathrm{O}$ とする。次の問いに答えなさい。}
\QQ{点 $\mathrm{B}$ の座標を $t$ を用いて表しなさい。}{}
\QQ{$x$ の値が点 $\mathrm{A}$ の $x$ 座標から点 $\mathrm{B}$ の $x$ 座標まで増加するとき、放物線における変化の割合を $t$ を用いて表しなさい。}{}
\QQ{点 $\mathrm{C}$ の座標を $t$ を用いて表しなさい。}{}
\QQ{$\Tri{OAB}$ の面積が $35$ であるとき、$t$ の値と点 $\mathrm{B}$ の座標を求めなさい。}{}
